\section{Discussion and Limitations}
A planner that trusts a source-calibrated region on unseen deployment data is, on this evidence, trusting a region
that is too small. At $\alpha=0.10$, the confirmatory models' cross-dataset gap is $+\gapRawConf$ points on nuScenes,
$+\gapRawConfRev$ points on AV2, and $+\wmGapAvW$ to $+\wmGapNsW$ points on Waymo, never zero. The gap is silent:
nothing in the deployed system signals it. That is why we treat measurement and audit as first-class contributions
rather than a prelude to the fix. Four results bear directly on practice.

Covariate reweighting is not a safe default. It widened the gap on every exploratory model we tried it on
(Sec.~\ref{sec:h2h3}), so a practitioner who reweights and stops is worse off than one who does nothing. The
label-free normalised score is not a substitute for labels either. It repaired the gap by 12--49\% on the exploratory
zoo and by 27\% on the confirmatory AV2 model, but made the confirmatory reverse-direction gap slightly worse. Nothing
observable at deployment time, without target labels, currently tells a practitioner which regime they are in.

The coverage audit (Sec.~\ref{sec:audit}) is the one method that behaved as pre-registered without exception, on every
model we tried it on: with on the order of $10^3$ labelled target scenes it detects a violation of this size reliably
and at a low false-alarm rate. The scene-level repair (Sec.~\ref{sec:scene}) is the one that turns a broken guarantee
into a working one, but its practical operating point is narrower than we first thought. Across three independent
confirmatory tests of the same pre-registered bar, on two datasets and two source models (Sec.~\ref{sec:confirm},
Sec.~\ref{sec:waymo}), the rule clears 88\% of draws reaching 90\% held-out coverage at 30--40 labelled scenes and
falls short at 50--60. We recommend the smaller budget, not the wider range the exploratory zoo appeared to support,
and we report the shortfall rather than the version of the story that would have looked cleaner.

The practical recommendation these four results support is concrete. Budget for 30--40 labelled target scenes,
collected as whole scenes rather than as a fixed count of agents, and recalibrate with the scene-level betting rule.
Treat covariate reweighting and the normalised score as diagnostics worth trying, not guarantees worth relying on.
Audit with a few hundred to a thousand labelled scenes before trusting any transferred region that has not been
recalibrated this way.

\paragraph{Does a stronger model fix it?} No. The confirmatory AV2 model is a properly resourced AutoBot, trained on
the full 183k-scene set and within 0.5\% of UniTraj's published accuracy, and its coverage gap (+15.1 points) is
larger than any model in the exploratory zoo (3.3 to 10.5 points). Four models now show the same pattern: point
accuracy and coverage validity are not the same property, and improving one does not repair the other. We do not have
a causal account of why the gap grows with accuracy. A larger model zoo, spanning architectures and training
recipes, would be needed to explain it, and we report the pattern as an open question rather than force an
explanation onto five data points.

\paragraph{Does it generalise beyond one dataset pair?} Largely yes. Waymo Open Motion (Sec.~\ref{sec:waymo}) was
never touched during method design: every hypothesis about it was fixed before any Waymo prediction was computed.
The direction of the failure, under-coverage, held for both source models and was larger than on nuScenes or AV2.
The repair transferred at the same budget that worked on nuScenes, for one of the two source models, and fell short
of the pre-registered bar for the other. We read this as qualified support: the failure mode generalises across three
independently collected datasets, and the fix generalises at the smaller label budget, but a practitioner should not
assume the 50--60-scene operating point without measuring it directly on their own data.

\paragraph{Limitations.}\label{sec:limitations}
(i) One architecture family (AutoBot). A second architecture, Wayformer, was planned and is supported by the
prediction and injection code, but training it within the time available for this submission would not have left
room for the pre-registered confirmatory and Waymo runs above, so it is left for future work rather than included
half-finished. (ii) The observed gap is an association. Its cause, whether annotation convention, sampling rate, or
agent mix, is not fully isolated: Sec.~\ref{sec:where} tests several candidates directly and finds that the measured
ones explain a minority of the gap, but unmeasured candidates cannot be ruled out. (iii) Coverage is marginal, not
conditional on any single scene; Sec.~\ref{sec:where}'s manoeuvre-level breakdown is a diagnostic, not itself a
guarantee. (iv) Hypotheses added after the first measurement pass are labelled post-hoc throughout and were
registered before the runs that tested them, but they were still informed by having seen Run~1. (v)
Proposition~\ref{prop:ident} covers scale shifts only; a more general shift could in principle be label-free
repairable in ways our injection experiments did not test for. (vi) The exploratory zoo's CPU-trained models remain
quantifiably below published state-of-the-art point accuracy, and every exploratory-only result (Sec.~\ref{sec:where}
and Sec.~\ref{sec:repair-scene}) should be read as designed-on rather than confirmed-on evidence until Sec.~\ref{sec:confirm}
and Sec.~\ref{sec:waymo} report the same quantity on a gated model, which they now do for the great majority of the
paper's claims.
